\section{Compact command groups}\label{sec:compact}
Theorem~\ref{thm:group-budget} applies to any orthogonal or unitary
representation, including noncommuting commands. The following two
consequences distinguish a finite command alphabet from a genuinely
compact command alphabet.

\begin{corollary}[Characters of compact groups]\label{cor:characters}
Let a finite command alphabet in a compact group contain the identity
and $g$, and let $\chi$ be a continuous one-dimensional unitary character
with $\chi(g)=e^{2\pi i\alpha}$, $\alpha$ irrational. Any wordwise
finite-register experiment whose terminal specification forces a
conditional $\chi$-amplitude at least $\kappa>0$ satisfies
\eqref{eq:occupation} and \eqref{eq:nonuniform} on these commands.
In particular, a badly approximable character angle forces
$\Omega(N^{1/3})$ hidden width.
\end{corollary}
\begin{proof}
Use packets $g^J e^{B-J}$. Their character images are exactly those in
Lemma~\ref{lem:separated}. The other commands may be assigned probability
zero by the test law. Apply Theorem~\ref{thm:group-budget} on $\C$.
\end{proof}
This includes higher-dimensional tori through their characters. It is a
lower bound for a selected observed character, not a multiplicative
claim about simultaneous independent outputs of several characters.
The next result uses a different orbit geometry rather than a character.

Fix an integer $d\ge2$ and a unit vector $v\in\R^d$. Consider the
separate experiment with commands $g_t\in\SO(d)$, supplied as exact
atomic matrices, and final coordinate query $j\in\{1,\ldots,d\}$.
Its specified mean is
\begin{equation}\label{eq:sphere-target}
          \E[B\mid g_1,\ldots,g_N,j]
             =\frac1{10}e_j^Tg_N\cdots g_1v.
\end{equation}
A realizing finite-state machine has Borel stochastic rows $T_{t,g}$.
Its width still counts only labels; access to this compact input and its
row function is part of this model, not an implementation by finite
binary input symbols. All terminal probabilities lie in $[9/20,11/20]$.

\begin{lemma}[Sphere quantization]\label{lem:sphere}
For normalized surface measure on $\mathbb S^{d-1}$ there are constants
$c_d,C_d>0$ such that the orbit distortion of Haar measure on $\SO(d)$
satisfies
\[
             \Gamma_k(\mathrm{Haar})\ge c_d k^{-2/(d-1)}.
\]
For every $0<\rho\le1/2$ there is a finite set $V\subset\mathbb S^{d-1}$
of at most $C_d\rho^{-(d-1)}$ points whose convex hull contains the
Euclidean ball of radius $\cos\rho$.
\end{lemma}
\begin{proof}
Haar rotation sends every unit vector to normalized surface measure.
A Euclidean radius-$r$ cap, $0<r\le1$, has angular radius
$2\arcsin(r/2)\le\pi r/2$. Polar integration bounds its probability
by $A_dr^{d-1}$ for some $A_d\ge1$: its integral is proportional to
$\int_0^{2\arcsin(r/2)}\sin^{d-2}u\,du$.
Take $r=(2A_dk)^{-1/(d-1)}$ in Proposition~\ref{prop:small-ball}.
The bound follows with $c_d=(2A_d)^{-2/(d-1)}/4$.

For the second assertion choose a maximal set with pairwise angular
distance greater than $\rho$. The angular radius-$\rho/2$ caps are
disjoint and each has surface mass at least $a_d\rho^{d-1}$, by the
same polar integral and $\sin u\ge2u/\pi$ on $[0,\pi/2]$.
Thus the set has at most $a_d^{-1}\rho^{-(d-1)}$ points.
Maximality makes it a $\rho$-net. Its support function in every unit
direction is at least $\cos\rho$, proving the ball containment.
These elementary estimates are the usual spherical quantization and
covering mechanisms; their role here is sequential compression.
\end{proof}

\begin{theorem}[Sharp width for all rotation commands]\label{thm:sphere}
For fixed $d\ge2$ and $0\le\varepsilon<1/(20\sqrt d)$, the least
hidden-state width for \eqref{eq:sphere-target}, in its compact-command
model, satisfies
\begin{equation}\label{eq:sphere-width}
             W^{\SO(d)}_{N,\varepsilon}=\Theta_{d,\varepsilon}
                                      (N^{(d-1)/2}).
\end{equation}
Writing $\kappa_d=1/(10\sqrt d)-2\varepsilon$, every profile obeys
\begin{equation}\label{eq:sphere-budget}
       c_d\sum_{t=1}^N K_t^{-2/(d-1)}\le\log(1/\kappa_d).
\end{equation}
The exact upper construction has common command rows and an
$N$-dependent terminal decoder.
\end{theorem}
\begin{proof}
Use independent Haar commands as the test law and set
$Y_t=g_t\cdots g_1v$. Combine the coordinate decoder means into a vector
$d(S_N)$ of norm at most $\sqrt d$. Wordwise approximation implies
\[
 \E\langle Y_N,d(S_N)\rangle\ge1/10-2\sqrt d\varepsilon,
 \qquad q_N\ge\kappa_d.
\]
Theorem~\ref{thm:group-budget} and Lemma~\ref{lem:sphere}, one command
per packet, give \eqref{eq:sphere-budget} and the lower bound.

For the upper bound take $\rho=1/(2\sqrt N)$ and a net $V$ as in
Lemma~\ref{lem:sphere}. Put $P=\conv V$, $\lambda=(\cos\rho)^{-1}$,
and $r_t=\lambda^t/4$. A label $u\in V$ has associated vector $r_tu$.
For any command $g$, the vector $gu/\lambda$ lies in
$(\cos\rho)\mathbb B^d\subset P$. Fix a Borel convex decomposition
of every point of $P$ into $V$, and use the coefficients of $gu/\lambda$
as the command row from label $u$. Such decompositions may be chosen
by a finite triangulation of $P$ using its vertices, with ties between
simplices resolved in a fixed order. The resulting rows are Borel,
normalized, and independent of $t$.

The expected next associated vector is exactly $g(r_tu)$.
Since $-\log\cos\rho\le\rho^2$ for $\rho\le1/2$,
$r_N\le e^{1/4}/4<1$. Also $v/10\in r_0P$, since
$\cos(1/2)/4>1/10$. Initialize by its chosen convex decomposition and
decode with means $r_Ne_j^Tu$. This realizes \eqref{eq:sphere-target}
exactly and uses at most $C_d(2\sqrt N)^{d-1}$ labels (and two output
labels, absorbed by the bound). No vectors are retained at runtime.
\end{proof}

For $d\ge3$ the commands here are noncommuting. For $d=2$ the order
$N^{1/2}$ concerns all circle rotations as inputs, while Theorem~\ref{thm:main}
has the smaller fixed alphabet $\{I,R_\alpha\}$ and order $N^{1/3}$.
These statements concern different command sets. The compact theorem
is not a claim of that exponent for an arbitrary finite generating set,
or an uncharged approximation of a real command by a finite word.
